% Reemplace las secciones de ejemplo con sus propias políticas y procedimientos.
\section{Propósito y alcance}
Este manual describe cómo la organización protege sus sistemas y sus datos.
Se aplica a todas las personas empleadas, contratistas y proveedores que tienen
acceso a los sistemas de la empresa.

\begin{policybox}
  Todos los datos de la empresa deben clasificarse, guardarse en un sistema
  aprobado y compartirse solo con quienes los necesitan para su trabajo.
\end{policybox}

\subsection{Roles y responsabilidades}
Cada rol tiene un responsable claro. La tabla indica quién decide, quién opera y
a quién hay que informar.

\begin{table}[H]
  \centering
  \begin{tabularx}{\linewidth}{@{}lXl@{}}
    \toprule
    \textbf{Rol} & \textbf{Responsabilidad} & \textbf{Contacto} \\
    \midrule
    Responsable de seguridad & Es dueño de este manual y aprueba las excepciones. & Equipo de seguridad \\
    Dueño del sistema & Mantiene al día las listas de acceso y los parches. & Equipo de plataforma \\
    Colaborador & Cumple las políticas e informa los incidentes. & Todas las personas \\
    \bottomrule
  \end{tabularx}
\end{table}

\section{Gestión de accesos}
El acceso se concede por rol y se revisa cada trimestre. No se permiten cuentas
compartidas; cada acción debe poder atribuirse a una sola persona.

\begin{procedurebox}
  \begin{steps}
    \item La jefatura presenta una solicitud de acceso con el sistema y el rol.
    \item El dueño del sistema verifica que el rol corresponde al cargo.
    \item El equipo de plataforma concede el acceso y anota el número del ticket.
    \item La persona confirma el acceso en un plazo de dos días hábiles.
  \end{steps}
\end{procedurebox}

\begin{warningbox}
  Nunca comparta credenciales por chat ni por correo. Si sospecha que una
  contraseña quedó expuesta, cámbiela primero e infórmelo después.
\end{warningbox}

\section{Plan de respuesta a incidentes}
Cuando algo parece estar mal, el primer objetivo es contener el problema y el
segundo es entenderlo. El flujo muestra el recorrido de todo incidente.

\begin{mermaid}[style=editorial,width=0.9\linewidth,
    caption={Flujo de respuesta a incidentes},pos=H]
  flowchart LR
  A[Detectar] --> B[Clasificar] --> C[Contener] --> D[Recuperar] --> E[Revisar]
\end{mermaid}

La persona de guardia puede listar los inicios de sesión fallidos recientes con
el comando siguiente.\par\vspace{2pt}

\begin{code}[lang=bash]
grep "authentication failure" /var/log/auth.log | tail -n 20
\end{code}

\begin{notebox}
  Anote en el ticket cada acción y su hora mientras trabaja. La revisión
  posterior depende de ese registro.
\end{notebox}

\section{Registro de decisiones}
Las decisiones importantes se registran para que quien lea después sepa por qué
se tomaron.

\begin{decision}{DR-004}{Adoptar llaves de seguridad físicas}
  \decisionpart{\lblContext}
  Los intentos de suplantación aumentaron en los dos últimos trimestres y los
  códigos de un solo uso todavía pueden interceptarse.

  \decisionpart{\lblDecision}
  Todo el personal con acceso a producción usará llaves de seguridad físicas
  para iniciar sesión desde el primero de enero.

  \decisionpart{\lblConsequences}
  Cada persona necesita dos llaves y se guarda una de repuesto en la caja fuerte
  de la oficina. Soporte debe atender la pérdida de una llave en un día hábil.
\end{decision}

\clearpage
\section*{\lblGlossary}
\addcontentsline{toc}{section}{\lblGlossary}
\begin{termlist}
  \item[Clasificación] La etiqueta que indica hasta dónde puede compartirse un documento.
  \item[Incidente] Todo evento que daña o puede dañar la confidencialidad, la integridad o la disponibilidad de los datos de la empresa.
  \item[Guía operativa] Una lista breve y probada de pasos para una tarea recurrente.
  \item[Mínimo privilegio] Dar a una persona solo el acceso que necesita para su trabajo actual.
  \item[Revisión] Una comprobación programada de que una política, un rol o una lista de acceso sigue siendo correcta.
  \item[Dueño del sistema] La persona responsable de un sistema y de su lista de acceso.
\end{termlist}
